\documentclass[12pt]{article}
\usepackage{amsmath, amssymb, amsthm}
\usepackage[margin=1in]{geometry}

\title{Project Euler Problem 368: A Kempner-like Series}
\author{}
\date{}

\begin{document}
\maketitle

\section{Problem Statement}
Compute the sum of the series
\[
S = \sum_{\substack{n=1 \\ \text{``506''} \not\subset n}}^{\infty} \frac{1}{n}
\]
where the sum runs over all positive integers whose decimal representation does not contain ``506'' as a substring. Report the answer to 10 decimal places.

\section{Solution}

\subsection{Convergence}
Removing all integers containing a fixed $k$-digit substring makes the harmonic-like series converge. Among $d$-digit numbers, the fraction avoiding a given 3-digit string is approximately $(1 - 10^{-3})^{d-2}$, which decays exponentially. Hence the tail contributions shrink geometrically and the series converges.

\subsection{Digit DP with Automaton}
We build an Aho-Corasick (or simply KMP) automaton for the pattern ``506'':
\begin{itemize}
\item \textbf{State 0}: No prefix of ``506'' matched.
\item \textbf{State 1}: Suffix ``5'' matched.
\item \textbf{State 2}: Suffix ``50'' matched.
\item \textbf{State 3}: Full match ``506'' (reject).
\end{itemize}

For each digit appended, we transition between states. We reject any number that reaches State~3.

\subsection{Computing the Sum}
Group integers by their digit length~$d$. For each $d$:
\begin{enumerate}
\item Use the digit DP with the automaton to enumerate $d$-digit numbers avoiding ``506''.
\item Compute $\sum 1/n$ over these numbers using subdivision:
    \begin{itemize}
    \item Split the range $[10^{d-1}, 10^d - 1]$ into blocks defined by the first few digits.
    \item Within each block, use the count of valid numbers and a smooth approximation (midpoint or Euler-Maclaurin) for $\sum 1/n$.
    \end{itemize}
\item For large $d$, the total contribution from $d$-digit numbers decays like $0.999^d$, allowing us to truncate.
\end{enumerate}

\subsection{Precision}
We compute contributions for $d = 1, 2, \ldots, D$ where $D$ is chosen so that the tail is below $10^{-12}$. With careful numerical summation, we achieve the required 10-digit precision.

\section{Answer}

\[
\boxed{253.6135092068}
\]

\end{document}
